\documentclass[conference]{IEEEtran}
\IEEEoverridecommandlockouts
\usepackage{cite}
\usepackage{amsmath,amssymb,amsfonts}
\usepackage{algorithmic}
\usepackage{graphicx}
\usepackage{textcomp}
\usepackage{xcolor}
\usepackage{hyperref}
\usepackage{orcidlink}
\usepackage[utf8]{inputenc}
\def\BibTeX{{\rm B\kern-.05em{\sc i\kern-.025em b}\kern-.08em
    T\kern-.1667em\lower.7ex\hbox{E}\kern-.125emX}}
\begin{document}

\title{Conversión de Señales mediante Circuitos con Amplificadores Operacionales}\\

\author{%
    \IEEEauthorblockN{Ana Freitez\, \orcidlink{0009-0009-6506-499X}}
    \IEEEauthorblockA{\textit{Facultad de Ingeniería} \\ 
    \textit{Universidad Latina de Panamá} \\ 
    Ciudad de Panamá, Panamá}
    \and
    \IEEEauthorblockN{Cristell Cedeño\, \orcidlink{0009-0002-2617-663X}}
    \IEEEauthorblockA{\textit{Facultad de Ingeniería} \\ 
    \textit{Universidad Latina de Panamá} \\ 
    Ciudad de Panamá, Panamá}
}

\maketitle

\begin{abstract}
% Resumen breve (150-250 palabras): objetivo del proyecto, circuito implementado
% (integrador y derivador en cascada con amplificadores operacionales),
% metodología utilizada (simulación y/o implementación física) y principales
% resultados obtenidos.
[Completar resumen del proyecto aquí.]
\end{abstract}

\begin{IEEEkeywords}
amplificador operacional, integrador, derivador, procesamiento de señales, circuitos analógicos
\end{IEEEkeywords}

\section{Introducción}
% Contexto general del tema, motivación del proyecto y objetivo principal.
% Mencionar brevemente qué problema resuelve un circuito integrador/derivador
% en cascada y por qué es relevante estudiarlo.
[Completar introducción.]

\subsection{Objetivos}
\begin{itemize}
    \item Objetivo general: [Completar]
    \item Objetivos específicos:
    \begin{itemize}
        \item [Completar]
        \item [Completar]
    \end{itemize}
\end{itemize}

\section{Marco Teórico}
% Explicar brevemente el principio de funcionamiento de un amplificador
% operacional en configuración integradora y derivadora, incluyendo las
% ecuaciones fundamentales.

\subsection{Amplificador Integrador}
[Completar explicación teórica del integrador. Incluir la ecuación característica.]
\begin{equation}
V_{out}(t) = -\frac{1}{RC}\int V_{in}(t)\,dt
\label{eq:integrador}
\end{equation}

\subsection{Amplificador Derivador}
[Completar explicación teórica del derivador. Incluir la ecuación característica.]
\begin{equation}
V_{out}(t) = -RC\,\frac{dV_{in}(t)}{dt}
\label{eq:derivador}
\end{equation}

\subsection{Consideraciones de Diseño Práctico}
% Mencionar por qué se añaden componentes adicionales (resistencia en paralelo
% al capacitor del integrador, resistencia en serie y capacitor en paralelo
% en el derivador) para evitar saturación y ruido en la implementación real.
[Completar.]

\section{Metodología}

\subsection{Diseño del Circuito}
% Describir el diseño propuesto: etapa integradora seguida de etapa
% derivadora en cascada. Referenciar la Fig.~\ref{fig:diagrama}.
[Completar descripción del diseño.]

\begin{figure}[htbp]
\centerline{\includegraphics[width=0.9\linewidth]{diagrama_circuito.png}}
\caption{Diagrama esquemático del circuito integrador-derivador en cascada.}
\label{fig:diagrama}
\end{figure}

\subsection{Materiales y Equipo}
% Completar la tabla con los componentes y equipos utilizados.
\begin{table}[htbp]
\caption{Lista de Materiales y Equipos}
\label{tab:materiales}
\centering
\begin{tabular}{l l c}
\hline
\textbf{Componente} & \textbf{Especificación} & \textbf{Cantidad} \\
\hline
Amplificador operacional & 741 & 2 \\
Resistencia & 10 k$\Omega$ & [Completar] \\
Resistencia & 100 k$\Omega$ & [Completar] \\
Capacitor & 0.1 $\mu$F & [Completar] \\
Protoboard & -- & 1 \\
Generador de señales & -- & 1 \\
Osciloscopio & -- & 1 \\
Fuente de alimentación & -- & [Completar] \\
\hline
\end{tabular}
\end{table}

\subsection{Procedimiento}
% Describir paso a paso cómo se realizó la simulación y/o el montaje físico.
\begin{enumerate}
    \item [Completar]
    \item [Completar]
    \item [Completar]
\end{enumerate}

\section{Resultados}

\subsection{Simulación}
% Insertar capturas del circuito simulado y del osciloscopio virtual.
[Completar descripción de los resultados de simulación.]

\begin{figure}[htbp]
\centerline{\includegraphics[width=0.9\linewidth]{simulacion_esquematico.png}}
\caption{Esquemático simulado en Proteus 8 Profesional.}
\label{fig:sim_esq}
\end{figure}

\begin{figure}[htbp]
\centerline{\includegraphics[width=0.9\linewidth]{simulacion_osciloscopio.png}}
\caption{Formas de onda obtenidas en el osciloscopio virtual: señal de entrada, salida del integrador y salida del derivador.}
\label{fig:sim_osc}
\end{figure}

\subsection{Implementación Física}
% Insertar fotografías del montaje en protoboard y capturas del osciloscopio real.
[Completar descripción del montaje físico.]

\begin{figure}[htbp]
\centerline{\includegraphics[width=0.9\linewidth]{montaje_protoboard.png}}
\caption{Montaje físico del circuito en protoboard.}
\label{fig:montaje}
\end{figure}

\begin{figure}[htbp]
\centerline{\includegraphics[width=0.9\linewidth]{osciloscopio_real.png}}
\caption{Capturas del osciloscopio real: (a) señal de entrada y salida del integrador; (b) salida del integrador y salida del derivador.}
\label{fig:osc_real}
\end{figure}

\subsection{Datos Medidos}
% Completar la tabla con los valores obtenidos en simulación y/o en el
% laboratorio para distintas frecuencias de entrada.
\begin{table}[htbp]
\caption{Amplitudes de Salida Medidas en Función de la Frecuencia de Entrada}
\label{tab:datos}
\centering
\begin{tabular}{c c c c}
\hline
\textbf{$f$ (Hz)} & \textbf{$V_{in}$ (V)} &
\textbf{$V_{out,int}$ (V)} & \textbf{$V_{out,der}$ (V)}\\
\hline
25    & 4 & [Completar] & [Completar] \\
50    & 4 & [Completar] & [Completar] \\
100   & 4 & [Completar] & [Completar] \\
200   & 4 & [Completar] & [Completar] \\
500   & 4 & [Completar] & [Completar] \\
1000  & 4 & [Completar] & [Completar] \\
\hline
\end{tabular}
\end{table}

\section{Análisis de Resultados}
% Comparar los resultados de simulación con los resultados físicos.
% Explicar cualquier diferencia observada (ruido, saturación, tolerancias
% de componentes, etc.) y relacionarlos con la teoría de la Sección II.
[Completar análisis comparando teoría, simulación e implementación física.]

\section{Conclusiones}
% Resumir los hallazgos principales y si se cumplieron los objetivos
% planteados en la introducción.
[Completar conclusiones.]

\section*{Agradecimientos}
[Completar, si aplica.]

\begin{thebibliography}{00}
\bibitem{b1} [Completar referencia 1.]
\bibitem{b2} [Completar referencia 2.]
\bibitem{b3} [Completar referencia 3.]
\end{thebibliography}

\end{document}